\documentclass[11pt,letterpaper]{article}
\usepackage[T1]{fontenc}
\usepackage[utf8]{inputenc}
\usepackage{libertinus}
\usepackage[margin=0.9in,headheight=15pt]{geometry}
\usepackage{microtype}
\usepackage{mathtools,amssymb,amsthm}
\usepackage{booktabs,longtable,array,tabularx}
\usepackage{xcolor}
\usepackage{enumitem}
\usepackage{listings}
\usepackage{fancyhdr}
\usepackage{hyperref}
\usepackage{bookmark}
\definecolor{ink}{HTML}{183A50}
\definecolor{muted}{HTML}{52606B}
\definecolor{codebg}{HTML}{F5F6F7}
\hypersetup{colorlinks=true,linkcolor=ink,citecolor=ink,urlcolor=ink,
 pdftitle={A Technical Audit of a Wolfram Language Radical-Denesting Prototype},
 pdfauthor={Technical review of the supplied source},pdfsubject={Correctness, branches, algorithms, implementation defects, and repairs}}
\pagestyle{fancy}
\fancyhf{}
\fancyhead[L]{\small Radical denesting: technical audit}
\fancyhead[R]{\small\thepage}
\renewcommand{\headrulewidth}{0.3pt}
\setlength{\parindent}{1.25em}
\setlength{\parskip}{0.25em}
\setlength{\emergencystretch}{2em}
\setcounter{tocdepth}{2}
\setlist{nosep,topsep=4pt}
\lstdefinestyle{wl}{language=Mathematica,basicstyle=\ttfamily\footnotesize,
 keywordstyle=\color{ink},commentstyle=\color{muted},stringstyle=\color{muted},
 backgroundcolor=\color{codebg},frame=single,rulecolor=\color{black!12},
 breaklines=true,breakatwhitespace=false,columns=fullflexible,keepspaces=true,
 showstringspaces=false,upquote=true,tabsize=2,xleftmargin=0.5em,xrightmargin=0.5em,
 aboveskip=0.8em,belowskip=0.8em}
\lstset{style=wl}
\newcommand{\code}[1]{\texttt{\detokenize{#1}}}
\newcommand{\src}[1]{\textup{[source, lines #1]}}
\newcommand{\Q}{\mathbb{Q}}
\newcommand{\C}{\mathbb{C}}
\newcommand{\Z}{\mathbb{Z}}
\DeclareMathOperator{\Arg}{Arg}
\DeclareMathOperator{\disc}{disc}
\DeclareMathOperator{\lcm}{lcm}
\DeclareMathOperator{\depthrad}{rdepth}
\newtheorem{proposition}{Proposition}[section]
\newtheorem{lemma}[proposition]{Lemma}
\newtheorem{corollary}[proposition]{Corollary}
\theoremstyle{definition}
\newtheorem{example}[proposition]{Example}
\newtheorem{definition}[proposition]{Definition}
\newenvironment{finding}[2]{\subsection{#1: #2}}{}
\begin{document}
\begin{titlepage}
\vspace*{0.6in}
{\color{ink}\LARGE\bfseries A Technical Audit of a Wolfram Language Radical-Denesting Prototype\par}
\vspace{0.35in}
{\Large\code{Strad}, \code{DenestRadicals3}, and \code{denest11}\par}
\vspace{0.2in}
{\large Algorithm reconstruction, mathematical justification,\par
 branch counterexamples, implementation defects, and a repair strategy\par}
\vspace{0.45in}
{\large September 5, 2026\par}
\vspace{0.4in}
\noindent\textbf{Material reviewed.} The single source file supplied with the request,
reproduced without modification as \code{source.wl} in the accompanying archive.
The file has 585 physical lines. All source-line references in this article refer
to that copy, not to a reformatted version.

\vspace{0.2in}
\noindent\textbf{Evidence.} The review combines complete static inspection of the supplied
source, mathematical derivations, current official Wolfram Language documentation,
and 26 successful independent checks in Python/SymPy. The supplied program was
\emph{not} executed in a Wolfram kernel in this session. Wolfram regression tests
are included for reproducing the predicted failure paths; they are not presented
as an executed test report.

\vspace{0.2in}
\noindent\textbf{Main conclusion.} The minimal-polynomial/GCD idea is mathematically useful,
but the present implementation is not a trustworthy general-purpose simplifier.
It can validate the wrong branch of the original number, its factorization layer
uses branch-unsafe transformations, and its search machinery contains independent
ordering and indexing defects. These are repairable without discarding the central
algebraic idea.
\vfill
\noindent{\small Source SHA-256:\par
\texttt{693a2825f1e2b8912e55371b90dc8ed115bc}\par
\texttt{33ee707dd3a4521eb22bcd2ccc93}}
\end{titlepage}

\begin{abstract}
The supplied Wolfram Language program attempts to denest exact algebraic radical
expressions by combining syntactic normalization with a multiplier search. For each
multiplier it computes a minimal polynomial, intersects that polynomial with a
binomial over an algebraic coefficient field, solves a lower-degree common factor,
and selects a candidate using a zero test. This article reconstructs that process,
proves the mathematical statements that actually justify it, and separates them
from the heuristic decisions that drive the search.

The most serious defect is a change of target: after replacing an original number
$\alpha$ by $\rho=\alpha^q$, the implementation validates candidates against the
principal value $\rho^{1/q}$ instead of against $\alpha$. Exact counterexamples
show how this reverses a sign even when every minimal-polynomial and GCD computation
is correct. Other findings include unrestricted radical factorization, a comparator
that ignores its second argument, a stale-length slice after prime pruning, unchecked
empty solution lists, missing multiplier validation, and an ineffective candidate
cap. The article also identifies sound components, including denominator
rationalization by a minimal polynomial, and gives a conservative candidate kernel,
regression cases, and a prioritized redesign that makes correctness independent of
search quality.
\end{abstract}

{\small\setlength{\parskip}{0pt}\tableofcontents}
\clearpage

\section{Scope, evidence, and the intended contract}
\label{sec:scope}

\subsection{What the attachment establishes}
The attachment contains definitions, comments, and diagnostic print statements, but
no usage messages, version declaration, sample invocation, test suite, specification
of supported inputs, or stated mathematical completeness theorem. Accordingly, the
review does not infer an author, a release history, or an intended domain beyond
what the code itself supports. In particular, the suffixes \code{3} and \code{11}
are names, not evidence of documented algorithm versions.

The public-looking entry point is
\begin{lstlisting}
Strad[expr, alllevels, func]
\end{lstlisting}
with the actual optional arguments defined by
\begin{lstlisting}
DenestRadicals3[expr_, alllevels_: False, func_: denest11]
\end{lstlisting}
The wrapper first performs custom factorization, then locates rational powers of
radicands for which \code{AlgebraicQ} succeeds. It can therefore occur in a symbolic
host expression, although the algebraic core is naturally suited to explicit,
finite, exact algebraic numbers. Those two domains are not equivalent, and the code
does not formally distinguish them.

A reasonable correctness contract would be:
\begin{quote}
For every supported input, return an expression with exactly the same value and
principal-branch interpretation. Replace it by a different representation only
when a declared simplification objective improves. If a bounded search does not
find such a representation, return the unchanged input with an explicit status.
\end{quote}
That is a proposed contract, not a claim that the attachment already implements it.

\subsection{Four levels of evidence}
Throughout the article, a \emph{source finding} is a consequence of the definitions
and control flow. A \emph{mathematical finding} is justified by a derivation given
here. An \emph{independent check} is an executed calculation in the accompanying
Python script. A \emph{Wolfram regression} is a supplied test whose actual kernel
result remains to be collected.

This distinction matters particularly for end-to-end examples. Exact identities and
polynomial GCDs can establish that an implemented branch choice is wrong without
establishing the precise printed transcript of a complete Wolfram evaluation.
Automatic simplification, factor ordering, and diagnostics can vary with kernel
version and expression representation. None of the predicted transcripts in this
article is described as an observed Wolfram run.

Official documentation is used to resolve language semantics, not to silently
replace the source algorithm. The core mathematical analysis and counterexamples
are derived from the supplied definitions. The companion files preserve this
separation: \code{verification_results.json} records only independent checks;
\code{regression_tests.wlt} contains unexecuted Wolfram assertions.

\subsection{Findings at a glance}
\small
\begin{longtable}{@{}p{0.065\linewidth}p{0.10\linewidth}p{0.56\linewidth}p{0.17\linewidth}@{}}
\toprule
ID & Priority & Finding & Source lines\\\midrule
\endfirsthead
\toprule ID & Priority & Finding & Source lines\\\midrule
\endhead
F1 & Critical & Candidates are certified against a reconstructed principal root rather than the original target. & 90--94, 122--127\\
F2 & Critical & The custom factorizer distributes and combines powers without the branch conditions needed for general inputs. & 527--543, 581--583\\
F3 & High & The prime comparator uses only its first operand, comparing a prime with its own exponent. & 177--181\\
F4 & High & Prime pruning computes a slice limit using the pre-pruning list length. & 184--194\\
F5 & High & Empty candidate lists are indexed at position 1; zero and other invalid multipliers are not rejected. & 107--134, 214--215\\
F6 & High & The default zero test is not an explicit exact-algebraic certificate; there is no final certificate against the original input. & 127, 314, 440\\
F7 & High & The nominal multiplier cap is not a cap; divisor and subset generation can allocate enormous lists. & 162--199, 402--420\\
F8 & Medium & A lower-degree GCD is called a denesting without measuring the returned representation. & 120--134, 309--324\\
F9 & Medium & Nested-power detection and the later rational-power algorithm have mismatched input contracts. & 84--108\\
F10 & Medium & Loading the source changes the global protection status of \code{Print}. & 1--3\\
F11 & Medium & Search heuristics discard possibilities and expose no completeness or exhaustion certificate. & 169--208, 224--248, 337--425\\
F12 & Low & Several controls are unused or misleading, notably \code{extrapower} and \code{maxprimes}. & 95--102, 162--165\\
\bottomrule
\end{longtable}
\normalsize
The priorities distinguish wrong answers from crashes, resource exhaustion, and
incomplete simplification. A heuristic denester is allowed to miss a simplification;
it is not allowed to replace a number by a different branch and call that success.

\section{Architecture and data flow}
\label{sec:architecture}

\subsection{The functions and their responsibilities}
\small
\begin{longtable}{@{}p{0.29\linewidth}p{0.57\linewidth}p{0.08\linewidth}@{}}
\toprule Function & Responsibility & Lines\\\midrule
\endfirsthead
\toprule Function & Responsibility & Lines\\\midrule
\endhead
\code{Strad} & Suppress diagnostic \code{Print} calls while invoking the wrapper. & 1--3\\
\code{DenestRadicals3} & Normalize an expression, mark radical subproblems, group some products, invoke a solver, and reassemble. & 4--76\\
\code{denest11} & Choose a reduction exponent and manage a search over multipliers and algebraic degrees. & 78--441\\
\code{trymultiplier} & Compute a minimal polynomial and a polynomial GCD, then generate and filter candidates. Local to \code{denest11}. & 107--155\\
\code{newmultipliers} & Generate new trial multipliers from a polynomial discriminant or from two trial records. Local to \code{denest11}. & 159--211\\
\code{RationalizeDenominator} & Express the inverse of an algebraic denominator using its minimal polynomial. & 443--451\\
\code{AlgebraicQ} & Ask whether an expression belongs to the algebraic numbers. & 453--454\\
\code{nestdepths} & Annotate an expression tree with a predicate-dependent nesting height. & 456--466\\
\code{nestdepthsort} & Flatten and group that annotation by height. & 468--480\\
\code{maxnestdepth} & Extract the portions retained at the maximum selected nesting height. & 482--498\\
\code{ComplexitySort} & Order expressions using leaf count, then absolute value. & 500--504\\
\code{Factorc} & Decode the custom nested factor representation into an expression. & 506--516\\
\code{fullfactor} & Construct and repeatedly refine that representation. & 518--551\\
\code{combine} & Extract common bases and exponents from factored summands. & 553--584\\
\bottomrule
\end{longtable}
\normalsize
The local \code{SumQ} predicate in the wrapper, lines 13--15, is a normalization
helper rather than another denesting algorithm.

\subsection{The top-level pipeline}
The wrapper's first operation is \code{Factorc[expr]}, before the algebraicity test
on individual radicands. It chooses \code{ReplaceAll} unless \code{alllevels} is
literally \code{True}; in that case it chooses \code{ReplaceRepeated}.
A matching rational power is temporarily represented as
\[
  f[\text{normalized radicand},\text{rational exponent}].
\]
The symbol \code{f} is a dynamically localized marker, not the mathematical function
being computed.

Normalization extracts a content factor, factors parts of that content, merges
factors recognized as sums or non-purely-imaginary complex atoms, and flips signs
of factors with numerically negative real part. It then multiplies the factors back
together, applies denominator rationalization, and simplifies.
The sign flips are accompanied by a compensating sign accumulator \code{j}; that
specific operation is not, by itself, a value-changing rewrite.

Next, some products of two marked radicals are bundled into a single one-argument
marker. The exact rule constructs
\[
 f[\mathrm{rad}_1^{p_1}\mathrm{rad}_2^{p_2}],
\]
not $f[(\mathrm{rad}_1\mathrm{rad}_2)^p]$. This is an important distinction: the
product-bundling rule itself does not assert an invalid product-of-roots identity.
The branch problem appears later when the bundled product is treated as the
principal root of its own power.

The remaining two-argument markers become \code{f[rad^pow]}. Another grouping rule
bundles products of marked values whose numerical real and imaginary parts are
both nonzero. The wrapper collects distinct \code{f[value]} occurrences and passes
their contents to \code{func}, defaulting to \code{denest11}.

\subsection{Single-level and all-level operation}
In the single-level mode, the collected markers are replaced with solver results
and the expression is simplified. The relevant limitation is the traversal and
matching behavior of the replacement rules; the name does not mean that exactly
one mathematical root is removed.

In all-level mode, the wrapper maintains a mixed list of unresolved markers and
resolved replacement rules. It scans for a marker containing no nested marker,
solves it, and propagates its replacement into the remaining unresolved entries.
This is a hand-written, bottom-up dependency resolution procedure. For an ordinary
finite, acyclic marked expression and a solver that returns marker-free values,
it has a natural progress measure: the number of unresolved markers decreases.
The code does not make that invariant explicit or check that a custom solver
respects it.

The final \code{ReplaceRepeated} reassembles the result after inner replacements
have made outer rule left-hand sides match. This explains why simply replacing it
with one \code{ReplaceAll} would not necessarily be a valid repair. A better
implementation would represent the dependencies explicitly or perform one controlled
bottom-up traversal.

\subsection{The core's state records}
An attempted multiplier is described by a triple
\[
  \{m,P_m,d_m\},
\]
where $P_m$ is stored as a pure function and $d_m$ is its degree.
The one-argument form of \code{MinimalPolynomial} really does return a pure-function
representation; subsequently calling \code{minpoly[arg1]} is valid, not a syntax
error~\cite{doc-minpoly}.

The local \code{trymultiplier} returns either
\begin{lstlisting}
{False, {multiplier, minpoly, degree}}
{True,  {multiplier, minpoly, degree}}
{True,  {multiplier, minpoly, degree}, answer}
\end{lstlisting}
The first Boolean means ``progress'' in the search, not necessarily ``a solution
was found.'' In particular, discovering a worse degree can count as progress when
it provides contrast with a previously seen degree.

A ready multiplier batch has the form
\begin{lstlisting}
{{doFullSet, originRecord}, {m1, m2, ...}}
\end{lstlisting}
and a deferred batch has the form
\begin{lstlisting}
{{False, smallerRecord}, newmultipliers,
 {smallerRecord, largerRecordOrNull, reductionpower}}
\end{lstlisting}
The scheduler recognizes a deferred batch because its second element is not a
list. At evaluation time it applies the stored function to the stored arguments.
A \code{history} list starts with the sentinel
\code{{0, Infinity, Infinity}} and is maintained in descending order of degree;
the last record therefore has the smallest recorded degree.

These structures are compact but overloaded. A tag or an association would make it
possible to validate records directly instead of inferring their meaning from list
lengths, positions, and whether their last element happens to be a list.

\subsection{What the nesting helpers measure}
For a Boolean predicate $w$ on expression nodes, define its weight to be one when
it is true and zero otherwise. The recursion in \code{nestdepths} computes
\[
 h(v)=w(v)+\max_{c\text{ an argument of }v}h(c),
\]
with the atom case handled separately. The result is an annotated tree of the form
\begin{lstlisting}
{expression, childRecord1, childRecord2, ..., height}
\end{lstlisting}
The usual predicate recognizes rational \code{Power} nodes. Thus the measure counts
matching nodes on an argument path; it is not ordinary Wolfram expression depth,
field-extension degree, or the minimum nesting depth among all equivalent formulas.

\code{nestdepthsort} rewrites the annotation into node/height pairs, groups equal
heights, and sorts the groups from largest to smallest. \code{maxnestdepth} instead
prunes the annotated tree relative to its maximum height and strips enclosing
metadata. Those utilities are used to select the reduction exponent and to compute
the diagnostic \code{extrapower}.

The representation retains the original expression at each annotated node while
also storing its children. In a long chain that is logically a repeated storage of
subexpressions; actual storage sharing is an implementation detail, not a complexity
bound supplied by this program. Because lists also encode the metadata, extending
the routines to arbitrary list-valued expressions or algebraic objects with internal
list structure requires care. A typed tree or a single traversal returning just the
required statistics would make the contract and tests substantially clearer.

\section{The algebraic idea, stated precisely}
\label{sec:algebra}

\subsection{Reduction exponent and trial numbers}
Let the original algebraic value be $\alpha$. The program chooses a positive integer
$q$ by taking an LCM of denominators extracted from selected maximum-depth rational
power nodes. It sets
\begin{equation}
  \rho=\alpha^q,\qquad \texttt{outerpower}=1/q.
  \label{eq:rho}
\end{equation}
The name ``outer root'' in the diagnostics should not be read as a theorem that this
$q$ is the denominator of every outer root in the input. The selection is based on
maximum predicate-dependent nesting height, and other branches may have different,
shallower rational powers.

For a trial multiplier $m$, the code constructs
\begin{equation}
  \beta_m=(m\rho)^{1/q},
  \qquad P_m(X)=\operatorname{MinPoly}_{\Q}(\beta_m).
  \label{eq:beta}
\end{equation}
All powers here have the principal-value meaning of Wolfram Language
\code{Power}~\cite{doc-power}. Suppose for now that $q\ge2$, $m\ne0$, and all
quantities are exact algebraic numbers.

Let $K$ be the algebraic coefficient field used by the GCD operation. With
\code{Extension -> Automatic}, Wolfram extends the coefficient field to include
algebraic numbers appearing in the polynomial arguments~\cite{doc-gcd}.
Consequently $m\rho\in K$. The field need not be the smallest field $\Q(m\rho)$:
its generators can depend on the explicit representation of the coefficients.

The critical computation is
\begin{equation}
  G_m(X)=\gcd_{K[X]}\bigl(P_m(X),X^q-m\rho\bigr).
  \label{eq:gcd}
\end{equation}

\begin{proposition}[Why the GCD is nonconstant]
Under the preceding hypotheses, $1\le\deg G_m\le q$.
\end{proposition}
\begin{proof}
The number $\beta_m$ is a root of both polynomials. Its minimal polynomial over $K$
therefore divides both in $K[X]$, and hence divides their GCD. The GCD has positive
degree and divides the degree-$q$ binomial, which gives the upper bound.
\end{proof}
Thus, a constant GCD in a valid, exact execution is not ordinary evidence that the
multiplier was unhelpful. It indicates a violated assumption, an unresolved
computation, or an implementation problem. The source prints a diagnostic in the
special case that the returned expression is literally \code{1}, but does not
propagate a structured failure.

\begin{proposition}[Candidate orbit]
If $r$ is any root of $G_m$, then
\begin{equation}
  \gamma_{r,j}=\frac{r}{m^{1/q}}\zeta_q^j,
  \qquad \zeta_q=e^{2\pi i/q},\quad 0\le j<q,
  \label{eq:orbit}
\end{equation}
satisfies $\gamma_{r,j}^q=\rho$. If $\rho\ne0$, this orbit contains every $q$th root
of $\rho$, including the original $\alpha$.
\end{proposition}
\begin{proof}
Because $G_m$ divides $X^q-m\rho$, one has $r^q=m\rho$. Also
$(m^{1/q})^q=m$, even for complex $m$, since the outer exponent is the integer $q$.
Taking the $q$th power of \eqref{eq:orbit} therefore gives $\rho$. If $\rho\ne0$,
the $q$ orbit elements are distinct; the binomial $X^q-\rho$ has exactly $q$
distinct roots, so the orbit contains all of them. For $\rho=0$ all candidates
collapse to zero, which is the only root.
\end{proof}
The roots-of-unity enumeration is therefore a sound way to handle the otherwise
unsafe division of principal roots. It does not require the generally false
identity $(m\rho)^{1/q}/m^{1/q}=\rho^{1/q}$.
What is essential is selecting the member that equals \emph{the original target}.

\subsection{What a smaller GCD degree does, and does not, prove}
If $e=\deg G_m<q$, solving $G_m(X)=0$ replaces a degree-$q$ equation over $K$ by a
lower-degree equation over $K$. That can produce a useful radical representation.
However, the coefficients of $G_m$ may themselves contain deeply nested radicals;
solving it may introduce more radicals or implicit \code{Root} objects. The
inequality $e<q$ is therefore a statement about a polynomial over a field, not a
proof that the displayed radical nesting depth has decreased.

Likewise, $d_m=\deg P_m$ is a degree over $\Q$, not the degree over $K$ and not a
measure of expression size. The program uses several different notions of
complexity without formally separating them:
\[
  d_m,\qquad e,\qquad \text{syntactic radical depth},\qquad
  \text{leaf count},\qquad \text{number of implicit algebraic objects}.
\]
They serve different purposes and cannot safely substitute for one another.

\section{A complete successful mathematical example}
\label{sec:worked}
Consider
\begin{equation}
  \alpha=\sqrt{5+2\sqrt6}=\sqrt2+\sqrt3.
  \label{eq:classic}
\end{equation}
The equality follows by squaring the positive right-hand side. Here $q=2$ and
$\rho=5+2\sqrt6$. For $m=1$,
\[
 P_1(X)=X^4-10X^2+1,
\]
since $(\alpha^2-5)^2=24$. Over $K=\Q(\sqrt6)$,
\[
 P_1(X)=(X^2-5-2\sqrt6)(X^2-5+2\sqrt6),
\]
so the GCD with $X^2-\rho$ has degree two. This trial gives no smaller equation.

For $m=2$,
\[
  \beta_2=\sqrt{10+4\sqrt6}=2+\sqrt6,
  \qquad P_2(X)=X^2-4X-2.
\]
The GCD becomes linear:
\[
  G_2(X)=X-2-\sqrt6.
\]
Dividing its root by $\sqrt2$ gives
\[
  \frac{2+\sqrt6}{\sqrt2}=\sqrt2+\sqrt3,
\]
which is exactly the desired denesting.

The following table gives independently checked legal trials. It is \emph{not} an
assertion that the unmodified scheduler visits them in this exact order.
\begin{center}
\small
\begin{tabular}{@{}cclc@{}}
\toprule
$m$ & $d_m$ & $P_m(X)$ & $\deg G_m$\\\midrule
$1$ & 4 & $X^4-10X^2+1$ & 2\\
$\sqrt6$ & 4 & $X^4-24X^2-6$ & 2\\
$2$ & 2 & $X^2-4X-2$ & 1\\
$3$ & 2 & $X^2-6X+3$ & 1\\
$\sqrt5$ & 8 & $X^8-490X^4+25$ & 2\\
\bottomrule
\end{tabular}
\end{center}
The discriminant of $P_1$ is
\[
 \disc(P_1)=147456=2^{14}3^2.
\]
This explains why searching multipliers formed from its discriminant primes can
find $2$ or $3$. It also illustrates why a multiplier can \emph{increase} the
minimal-polynomial degree: multiplying by $\sqrt5$ produces a degree-eight trial.

A useful distinction emerges. The original and denested expressions in
\eqref{eq:classic} represent the same degree-four number over $\Q$; denesting did
not reduce its intrinsic algebraic degree. The multiplier temporarily moved the
search to a number of degree two, and dividing out the multiplier restored a
simpler representation of the original degree-four number.

\section{How the multiplier search works}
\label{sec:search}

\subsection{Initial guesses from radicand terms}
Without a supplied multiplier list, the code expands a rationalized radicand,
attempts to remove rational coefficients from its terms, and derives complementary
radical factors. For a factor $v=b^{a/b_0}$, the exponent used in the complementary
construction is
\[
  \frac{b_0-a}{a},
\]
which formally produces $b^{(b_0-a)/b_0}$. This is a heuristic way to complete a
fractional exponent to an integer. The source prepends the multiplier $1$.

Several qualifications are important. The extraction uses
\code{List @@ Expand[...]} as though it always obtained the summands of a
\code{Plus}. But \code{Apply} replaces a head: on a product it returns factors,
on a power it returns the base and exponent, and on an atom it effectively does
nothing~\cite{doc-apply}. Thus the intended term interpretation is shape-dependent.
The expression should first be tested for head \code{Plus} and otherwise treated
as one term.

The removal of coefficients and the recombination of powers are proposal-generation
heuristics, not established equivalences between multipliers. They can be acceptable
in an incomplete search, provided that every final candidate is independently
certified against the original target.

\subsection{Degree-gated GCD computation}
For each multiplier, the minimal polynomial is always computed first. A GCD is
computed only when
\[
  d_m<d_{\mathrm{best}}
  \quad\text{or}\quad
  \gcd(d_m,d_{\mathrm{best}})<d_{\mathrm{best}}.
\]
The first comparison is useful as a heuristic because a rational-degree reduction
can signal an exploitable relation. The second is an attempt to exploit
incomparable reductions in degree.

Neither comparison is a completeness theorem. In particular, an equal-degree
trial is skipped after the best degree is known, although its coefficient-field
GCD and its resulting expression need not be identical to those of previous trials.
The program supplies no argument that all useful denestings must pass this gate.
A redesigned search should treat this comparison as a priority rule, not confuse it
with a logically conclusive rejection of a multiplier.

Nor does the numerical GCD of two field degrees establish a common field of that
degree. For example, $\Q(2^{1/4})$ and $\Q(3^{1/6})$ have degrees four and six.
Their intersection can only have degree one or two. The first field has the unique
quadratic subfield $\Q(\sqrt2)$ and the second the unique quadratic subfield
$\Q(\sqrt3)$, so their intersection is $\Q$, not a degree-two field. Uniqueness here
follows because two distinct quadratic subfields would generate a degree-four
biquadratic field: that is impossible in a degree-six field, and would make the
real pure-quartic field Galois, which it is not. This example is a warning about
degree-only inference, not an asserted counterexample to a particular generated
multiplier in the attachment.

\subsection{Discriminant-based batches}
In the usual branch of \code{newmultipliers}, the code partially factors
$\disc(P_m)$, retains factors passing \code{PrimeQ}, sorts them, trims certain large
prime gaps, and takes a prefix. If the chosen distinct primes are
$p_1,\ldots,p_k$, it forms
\begin{equation}
 \left\{\prod_{i=1}^k p_i^{e_i}:0\le e_i\le q-1\right\}.
 \label{eq:divisors}
\end{equation}
The implementation obtains this set by calling
\code{Divisors[(Times @@ multiplierprimes)^(q - 1)]}.
It unions the result with \emph{all} retained discriminant primes, drops the unit,
and multiplies the batch by the earlier successful or promising multiplier.

The union detail matters: a prime excluded from the selected prefix can still be
tried singly. It is excluded from the multi-prime divisor combinations, not
necessarily from every subsequent trial.

\code{FactorInteger[n, Automatic]} is a documented partial-factorization mode,
not an invalid call. It extracts factors that are easy to find and can leave a
composite residual~\cite{doc-factorinteger}. The source explicitly deletes that
residual instead of postponing it. Consequently, the strategy can lose the primes
inside a hard residual.

Discriminants are reasonable sources of candidates, but the code proves no theorem
that these candidates suffice. A polynomial discriminant can also include factors
arising from a chosen generating element rather than only intrinsic field
ramification. For a monic minimal polynomial of an algebraic integer, the familiar
relation is
\[
 \disc(P)=\disc(K)\,[\mathcal O_K:\Z[\beta]]^2.
\]
To see the squared index factor, write the power basis in an integral basis with
change-of-basis matrix $A$. The trace-pairing matrix transforms as $A^{\mathsf T}TA$,
so its determinant is multiplied by $(\det A)^2$, and $|\det A|$ is the index.
The source's general minimal polynomials need not be monic, so even this specific
interpretation requires hypotheses not imposed by the implementation.

\subsection{Batches from incomparable degree records}
When neither observed degree divides the other, \code{newmultipliers} instead
proposes the product of the two stored multipliers. It prints a polynomial GCD of
their rational minimal polynomials as a diagnostic. Distinct irreducible rational
minimal polynomials usually have GCD one; that diagnostic does not itself construct
a common algebraic field or justify the proposed product. The product is another
reasonable guess, not an established deduction.

\subsection{History, scheduling, and stopping}
The scheduler can split a partially processed batch and insert a newly generated
batch immediately after it, or defer it according to degree and leaf-count
comparisons. The initial full batches are allowed to finish so that their degree
statistics can be used to build additional combinations.

After an apparent success, the source may stop immediately if this was effectively
the first recorded trial, or if a ratio involving \code{history[[2,-1]]} and the
new degree is at least $q$. Because history is sorted by degree,
\code{history[[2,-1]]} is the largest finite degree currently recorded, not an
immutable baseline degree of the input. An unrelated worse trial can therefore
change this heuristic stopping criterion.

The variable \code{solution} is replaced by the latest accepted answer without
comparing it with the preceding solution under an explicit objective. Returning
the original expression at the end means only that this search did not retain an
answer. It is not a proof that no denesting exists.

\section{Critical defect F1: the wrong branch is certified}
\label{sec:branch}

\subsection{The precise change of target}
The decisive sequence is
\begin{lstlisting}
radicand = problem^reductionpower;
...
Select[possibilities,
       PossibleZeroQ[# - radicand^outerpower] &]
\end{lstlisting}
with \code{outerpower = 1/reductionpower}.
Thus the original target $\alpha$ is silently replaced by
\begin{equation}
 \widetilde\alpha=(\alpha^q)^{1/q}.
 \label{eq:wrongtarget}
\end{equation}
Taking the integer power in \eqref{eq:rho} is legitimate. Inverting that operation
with a principal root, without retaining a branch label, is not.

\begin{proposition}[Exactly when reconstruction preserves the value]
For $\alpha\ne0$ and an integer $q\ge2$, write
$\alpha=re^{i\theta}$ with $\theta\in(-\pi,\pi]$. Then
\begin{equation}
 (\alpha^q)^{1/q}=\alpha e^{-2\pi i k/q},
 \quad q\theta-2\pi k\in(-\pi,\pi].
 \label{eq:phase}
\end{equation}
In particular, the reconstruction equals $\alpha$ exactly when
$\theta\in(-\pi/q,\pi/q]$.
\end{proposition}
\begin{proof}
The principal argument of $\alpha^q$ is $q\theta-2\pi k$ for the unique integer $k$
placing it in $(-\pi,\pi]$. Taking its principal $q$th root gives
$r e^{i\theta-2\pi i k/q}$. In the allowed range of $\theta$, equality with $\alpha$
occurs precisely when $k=0$, yielding the stated sector. The asymmetry at the two
endpoints is inherited from the chosen principal-argument interval.
\end{proof}
A single principal $q$th root lies in this sector, which explains why many
ordinary positive-real examples work. A negative multiple, a product of such roots,
or a rational power with numerator greater than one need not lie in it.

\subsection{A direct negative-target counterexample}
Consider the core call
\begin{lstlisting}
denest11[-Sqrt[5 + 2 Sqrt[6]], {2}]
\end{lstlisting}
The nested radical is detected, $q=2$, and the radicand becomes $5+2\sqrt6$.
The multiplier $2$ produces exactly the linear GCD from Section~\ref{sec:worked}.
The candidate filter keeps $\sqrt2+\sqrt3$, since that equals the principal square
root of the radicand. But the original target is its negative.

This counterexample concerns \code{denest11} as defined. It should not be
misrepresented as saying that \code{Strad[-Sqrt[...]]} must lose the same sign:
the outer wrapper may leave a simple minus sign outside its marked radical and
thereby preserve it. The next two examples explain why the defect nevertheless
matters to the wrapper's supported-looking inputs.

\subsection{A single rational power with numerator three}
Let
\[
 z=-1+2i\sqrt2,
 \qquad u=1+i\sqrt2.
\]
Since $u^2=z$ and $u$ lies in the right half-plane, $u$ is the principal square root
of $z$. Therefore
\begin{equation}
 \alpha=z^{3/2}=u^3=-5+i\sqrt2.
 \label{eq:threehalf}
\end{equation}
The principal argument of $z$ is between $\pi/2$ and $2\pi/3$, so this is also
consistent directly with the principal logarithm definition of $z^{3/2}$.
Squaring gives
\[
 \rho=\alpha^2=23-10i\sqrt2.
\]
Its principal square root is $5-i\sqrt2=-\alpha$, not $\alpha$.

With $m=1$, the source's trial number is $\beta=5-i\sqrt2$, whose rational minimal
polynomial is
\[
 P(X)=X^2-10X+27.
\]
Over a coefficient field containing $i\sqrt2$,
\[
 \gcd(P(X),X^2-\rho)=X-5+i\sqrt2.
\]
Thus the very first forced trial can produce a linear GCD and select the wrong
branch. The corresponding core regression is
\begin{lstlisting}
denest11[(-1 + 2 I Sqrt[2])^(3/2), {1}]
\end{lstlisting}
The wrapper pattern accepts any rational exponent, not merely a numerator-one
root, so this is not an exotic exponent outside its pattern syntax.

\subsection{A bundled product of principal square roots}
Define
\[
 z_1=-1+2i\sqrt2=(1+i\sqrt2)^2,
 \qquad z_2=-2+2i\sqrt3=(1+i\sqrt3)^2.
\]
Both displayed square roots are principal. Their product is
\begin{equation}
 \alpha=\sqrt{z_1}\sqrt{z_2}
       =1-\sqrt6+i(\sqrt2+\sqrt3),
 \label{eq:product}
\end{equation}
whose real part is negative and imaginary part positive. The wrapper's two-marker
product rule can send this whole product to the core, because both radicands
contain powers.

For $q=2$, the reconstructed principal root is again $-\alpha$. The independently
checked minimal polynomial of $-\alpha$ is
\[
 P(X)=X^4+4X^3+4X^2+48X+144,
\]
and its GCD with $X^2-\alpha^2$ over the explicit algebraic coefficient field is
$X+\alpha$. The polynomial step is correct; it is solving for the wrong selected
branch.

An integration regression that isolates the first multiplier is
\begin{lstlisting}
DenestRadicals3[
 Sqrt[-1 + 2 I Sqrt[2]] Sqrt[-2 + 2 I Sqrt[3]],
 False,
 Function[t, denest11[t, {1}]]
]
\end{lstlisting}
The exact final representation of an actual run depends on the preprocessing
path, so this test is supplied as an unexecuted integration regression. The
algebraic failure of the reconstructed target is independently established.

\subsection{The minimal repair}
The original target must be an explicit argument to the candidate-selection step:
\begin{lstlisting}
Select[possibilities, ExactAlgebraicEqualQ[#, original] &]
\end{lstlisting}
The existing roots-of-unity orbit already contains the original target under the
hypotheses of Section~\ref{sec:algebra}. There is no need to guess a phase from
floating-point arguments. A final exact check of the fully reconstructed output
against the original input should be a second, independent boundary condition.

\section{Critical defect F2: branch-unsafe custom factorization}
\label{sec:factor}

\subsection{The custom representation}
\code{fullfactor} begins with \code{FactorList}, then converts entries into nested
lists representing products, powers, and sums. A typical leaf factor has the form
\code{{base, exponent}} at one stage; deeper lists encode factored summands.
\code{Factorc} reconstructs expressions through repeated patterns that turn
pairs into powers, collections into products, and a different nesting shape into
sums.

The representation has no distinct heads such as \code{ProductNode},
\code{PowerNode}, and \code{SumNode}. Instead, meaning is inferred from list depth
and cardinality. This makes the correctness of its transformations difficult to
state, and allows expression data and structural metadata to look alike.

\subsection{Two transformations requiring missing hypotheses}
The power case in \code{fullfactor} multiplies the exponent of a power base by the
exponent stored in its factor record. Algebraically, it reconstructs
\[
  (u^a)^b\longmapsto u^{ab}.
\]
The common-factor extraction in \code{combine} moves common bases outside an outer
power. If a sum has been written $F S$, the decoded result has the form
\[
  (FS)^b\longmapsto F^b S^b.
\]
Both are valid under appropriate restrictions, but not for arbitrary principal
complex powers. Wolfram's own \code{Power} only performs the corresponding general
automatic transformations when the outer exponent is an integer; unrestricted
power expansion requires additional assumptions~\cite{doc-power,doc-powerexpand}.

Some initial \code{FactorList} exponents are integers, so not every application of
the first rule is wrong. The problem is that repeated unwrapping and common-factor
extraction also operate with rational outer exponents and do not preserve the
branch conditions required at that point.

\subsection{A symbolic counterexample to common-factor extraction}
The factor records for $-x$ and $-y$ share the factor $-1$. When \code{combine}
extracts it from an outer square root, its reconstruction is
\begin{equation}
 \sqrt{-x-y}\longmapsto i\sqrt{x+y}.
 \label{eq:factorbad}
\end{equation}
At $x=y=-1$, the left side is $\sqrt2$, while the right side is $-\sqrt2$.
Thus the transformation encoded by the factor extraction cannot be an
unconditional identity.

The corresponding Wolfram regression intentionally performs substitution
\emph{after} symbolic preprocessing:
\begin{lstlisting}
original = Sqrt[-x - y];
changed = Factorc[original];
RootReduce[(changed - original) /. {x -> -1, y -> -1}]
\end{lstlisting}
A correct unconditional normalizer must produce zero. Substituting $x=y=-1$
\emph{before} calling \code{Factorc} would instead feed it the already evaluated
number $\sqrt2$ and would not test the same transformation. The exact wrapper
output of this regression is not claimed as an observed kernel result; the
branch-unsafety of \eqref{eq:factorbad} is a proved counterexample to the rule.

\subsection{Why the wrapper's algebraicity test does not protect it}
The first call is \code{Factorc[expr]}, outside and before the replacement pattern
\begin{lstlisting}
radicand_?AlgebraicQ^pow_Rational
\end{lstlisting}
Even a symbolic power that is subsequently
excluded from algebraic denesting has already passed through custom factorization.
Consequently, the wrapper cannot rely on the later pattern test to justify
unconditional transformations of the whole input.

Moreover, \code{ComplexExpand} assumes variables real unless they are specified
otherwise~\cite{doc-complexexpand}. That default is harmless for suitable explicit
numeric algebraic inputs, but it is not a replacement for a documented domain
contract when symbolic values reach this layer.

\subsection{A safer approach}
Preserve each principal \code{Power} node until a value-preserving transformation
has been certified. Integer outer powers may be distributed freely. Positive-real
common factors can be extracted from rational powers, but arbitrary signs and
complex phases cannot. A branch-aware expression tree should retain the original
power node rather than flattening it into a base/exponent record whose history has
been lost.

The simplest initial repair is to remove \code{Factorc} from the correctness-critical
path and use only transformations whose applicability is established for the exact
input. Search quality may temporarily decrease, but a correct unchanged expression
is better than a branch-changed ``simplification.''

\section{Definite implementation defects and unsafe failure paths}
\label{sec:implementation}

\begin{finding}{F3}{The prime comparator ignores its second operand}
The source contains
\begin{lstlisting}
#1[[1]] <= #[[2]] &
\end{lstlisting}
Here \code{#} is exactly \code{#1}, not \code{#2}~\cite{doc-slot}. The comparator is
therefore equivalent to
\begin{lstlisting}
Function[{left, right}, left[[1]] <= left[[2]]]
\end{lstlisting}
For records \code{{2,1},{3,1}}, it returns \code{False} in both directions,
because $2\nleq1$ and $3\nleq1$. It compares each prime with its own exponent, not
with the other prime. Prime-gap trimming and prefix selection downstream therefore
operate without the intended numerical ordering invariant.

Replace it with
\begin{lstlisting}
SortBy[factors, First]
\end{lstlisting}
or, minimally,
\begin{lstlisting}
Sort[factors, #1[[1]] <= #2[[1]] &]
\end{lstlisting}
Using \code{<=} is not the defect: Wolfram permits such non-strict ordering
functions, and its documentation explicitly discusses \code{GreaterEqual} for
preserving the relative order of equal elements~\cite{doc-sort}. Importing a C++
strict-weak-ordering objection here would miss the actual problem.
\end{finding}

\begin{finding}{F4}{A stale length is used after deleting pruned primes}
The source assigns zeros to selected prime entries and then evaluates
\begin{lstlisting}
multiplierprimes = DeleteCases[multiplierprimes, 0][[
 1 ;; Min[desiredCount, Length[multiplierprimes]]]];
\end{lstlisting}
The first part of the right-hand side has a shorter list, but the \code{Length}
inside the span still refers to the old value of \code{multiplierprimes}. The
assignment has not happened yet.

For a concrete source-invariant fixture, take the sorted primes
\[
 \{2,3,5,7,11,2003\},\qquad q=2.
\]
The last ratio is greater than 100, so the loop replaces $2003$ by zero and stops
at its minimum of five entries. The deleted list has length five. However,
\code{desiredCount} is ten and the old length is six, so the code requests the
span \code{1 ;; 6} of a five-element list. This is an out-of-range part request.
The fixture consists entirely of primes; it does not rely on a composite leftover.

The repair is to make the intermediate state explicit:
\begin{lstlisting}
multiplierprimes = DeleteCases[multiplierprimes, 0];
count = Min[desiredCount, Length[multiplierprimes]];
multiplierprimes = Take[multiplierprimes, count];
\end{lstlisting}
The faulty comparator can affect whether a particular real input reaches this
fixture. It does not invalidate the independent indexing defect; fixing the sort
alone is insufficient.
\end{finding}

\begin{finding}{F5}{Empty candidate sets and invalid multipliers are unchecked}
After solving and filtering, the source tests only whether the result has more
than one element. It then unconditionally evaluates \code{result[[1]]}.
An empty list therefore enters a \code{Part} failure path and can subsequently be
packaged as if it were an answer. The outer scheduler further identifies an answer
by checking whether the last element of the attempt record is not a list, which
is no substitute for validating a successful result type.

The forced list is accepted without checking its elements. In particular, with
$m=0$, $q=2$, and a nonzero target, the trial minimal polynomial is $X$ and the
GCD with $X^2$ is $X$. The code then tries to form $0/0$ while dividing the root by
$m^{1/q}$. Candidate filtering cannot repair that invalid arithmetic. A direct
failure regression is
\begin{lstlisting}
denest11[Sqrt[5 + 2 Sqrt[6]], {0}]
\end{lstlisting}
Other invalid inputs include inexact, symbolic, nonalgebraic, and nonfinite
multipliers. They can leave \code{MinimalPolynomial}, \code{PolynomialGCD}, or
\code{Solve} unevaluated, after which positional indexing propagates malformed
state.

Validate exact algebraicity, finiteness, nonzero multipliers, and an integer
$q\ge2$ before the trial. Check the return types and polynomial degrees of each
algebraic operation. If no certified candidate remains, return a distinct
\code{NoCandidate} status rather than indexing the list.
\end{finding}

\begin{finding}{F6}{The verification boundary is insufficient}
The default \code{PossibleZeroQ} is documented as a quick symbolic/numerical test
that is not always accurate. Current documentation provides
\code{Method -> "ExactAlgebraics"} for guaranteed methods on explicit algebraic
numbers~\cite{doc-possiblezero}. The source does not request that method, does not
first establish a precise exact-input contract, and---more decisively---passes the
wrong target to the test.

This is not a claim that a specific valid exact candidate is known to be
misclassified by the default method in this attachment. The demonstrable defect is
that the implementation lacks an explicit exact certification boundary, while F1
shows that even a perfect zero test is being asked the wrong question.

For explicit exact algebraic values, a conservative equality predicate is
\begin{lstlisting}
ExactAlgebraicEqualQ[a_, b_] :=
 Quiet[Check[RootReduce[a - b] === 0, False]]
\end{lstlisting}
provided the inputs have first passed an exact-algebraic validator.
\code{RootReduce} is designed to reduce algebraic combinations to a single exact
algebraic representation~\cite{doc-rootreduce}. An unresolved result is a failure
to certify, not permission to accept a numerical near-zero.
\end{finding}

\begin{finding}{F9}{The input gate and the subsequent algorithm disagree}
The initial nested-power detector matches arbitrary \code{Power} expressions,
including integer and symbolic exponents. The later reduction logic selects only
\code{Power[_, _Rational]}, and \code{trymultiplier} is defined only when the outer
power has the form \code{Rational[1, _Integer]}.
Thus passing the first detector does not imply that the reduction exponent is
well formed or that the local trial function will evaluate.

The inner \code{MemberQ[..., Infinity]} also excludes the base expression itself:
\code{Infinity} means levels 1 through infinity, whereas \code{{0, Infinity}}
includes the root~\cite{doc-memberq}. A power directly used as another power's
base is therefore not tested in the same way as a power inside a sum in that base.
Automatic evaluation may simplify many such cases first, so this is a structural
coverage issue, not a claim that every nested-power input is missed.

The helper \code{AlgebraicQ} is a domain-membership query, not an exactness
validator. Official documentation describes immediate recognition of explicit
algebraic constructions and permits unresolved membership questions for other
expressions~\cite{doc-algebraics}. The source adds no independent exact-precision
or finite-value check.

The repair is one explicit eligibility predicate, shared by detection, depth
measurement, reduction selection, and the solver. Purely symbolic expressions
should either be excluded from the numeric core or handled in a separate,
assumption-aware interface.
\end{finding}

\begin{finding}{F10}{Loading the file changes global kernel state}
The top of the source executes \code{Unprotect[Print]} and never restores its
previous protection status. The \code{Block} in \code{Strad} localizes the value
used for \code{Print}; it does not undo that top-level mutation of attributes.
This changes the surrounding kernel merely by loading the file.

More generally, \code{Block} uses dynamic scoping, while \code{Module} introduces
fresh symbols~\cite{doc-block}. The many dynamically localized names are not
automatically bugs: the source does localize the parent loop variables, including
\code{i3}. But private syntax markers, scratch variables, and callbacks are easier
to isolate using a package context and fresh lexical symbols. Deliberate dynamic
settings such as assumptions or logging configuration are separate concerns.

Replace global print suppression with an explicit logging function or a
\code{Verbose} option. Do not unprotect a system symbol as a permanent load-time
side effect. A package should also avoid \code{ClearAll} on generic global names
such as \code{combine} that may already belong to the caller's session.
\end{finding}

\section{Search limitations, resource risks, and maintainability}
\label{sec:resources}

\subsection{F7: the multiplier cap can be exceeded by orders of magnitude}
With $k$ distinct selected primes, the divisor set in \eqref{eq:divisors} has exactly
$q^k$ elements: each of the $k$ exponents independently has $q$ choices.
The source chooses a prefix length resembling
\[
 k=\min\!\left(L,\max\!\left(5,
       \left\lceil\frac{\log1000}{\log q}\right\rceil\right)\right),
\]
where $L$ is the available prime count, subject to the separate stale-length bug.
For sufficiently many available primes, the number of nonunit divisor candidates
is therefore:
\begin{center}
\begin{tabular}{@{}rrr@{}}
\toprule $q$ & Selected $k$ & $q^k-1$\\\midrule
2 & 10 & 1,023\\
3 & 7 & 2,186\\
10 & 5 & 99,999\\
100 & 5 & 9,999,999,999\\
\bottomrule
\end{tabular}
\end{center}
These counts omit extra singleton primes introduced by the union. The constant
\code{multipliercap = 1000} is therefore neither a strict bound on trials nor a
bound on allocated memory. The forced minimum of five primes is particularly
hazardous for high-order roots.

To bound the complete divisor product alone by $M$ nonunit candidates, one needs
$q^k-1\le M$, hence
\[
 k\le\left\lfloor\log_q(M+1)\right\rfloor.
\]
Even that correction is insufficient if another union adds more candidates, or if
subsequent batches have no global budget. A robust scheduler must impose a total
trial limit and generate candidates lazily rather than allocating the entire
cartesian product first.

\subsection{Exponential subset expansion after the initial pass}
The post-pass construction uses each selected multiplier to the powers 0 and 1
and distributes across the choices. For $r$ selected multipliers, this creates up
to $2^r$ combinations before duplicate elimination and simplification. Removing
duplicates afterward does not prevent the initial allocation.

Its simplifications strip integer factors and reduce some exponents modulo their
denominators. These operations are only heuristics on the proposed multiplier
set. In general, multiplying a multiplier by an arbitrary integer is not an
algebraically irrelevant change: Section~\ref{sec:worked} shows that $m=1$ and
$m=2$ have different trial degrees. Normalization modulo certified $q$th powers
would require its own justification and branch handling.

\subsection{No global cache, budget, or explicit exhaustion state}
There is no global set of already attempted exact multiplier values, no minimal-
polynomial cache, no total candidate budget, no degree or coefficient-size limit,
and no controlled time or memory budget around expensive algebraic operations.
The same multiplier can arise through different discriminant batches or products.
Syntactic \code{DeleteDuplicates} within one list is not an algebraic cache across
the search.

The lack of these controls does not, by itself, prove that a particular valid input
will loop forever. It does mean that the program offers no application-level bound
on work, no explicit proof of termination of its changing work queue, and no clean
distinction between ``budget exhausted'' and ``no candidate exists.'' A user-facing
implementation should expose those distinctions directly.

\subsection{Degree and coefficient growth}
Computing a minimal polynomial of a radical combination can require a compositum
of several algebraic extensions. Both its degree and the size of its coefficients
can grow rapidly as radicals and multipliers are combined. The following GCD may
work over another large explicit coefficient field, and discriminant factorization
adds integer-arithmetic costs unrelated to the visual size of the original input.

A useful cost model is therefore not just the number of multiplier trials, but a
sum of costs depending on rational degree, coefficient-field degree, coefficient
height, root order, and candidate representation size. The attachment provides no
measurements from which reliable wall-clock estimates can be inferred. In
particular, the number \code{1000} should not be interpreted as an experimentally
validated performance bound.

\subsection{F8: success is not measured against a denesting objective}
The comment ``progress because a denesting occurred'' follows only a smaller GCD
degree and a surviving equality test. There is no comparison of the answer's
radical depth with the input's, no prohibition on introducing \code{Root}, and no
comparison with a previously retained answer.

An explicit objective could first exclude new implicit algebraic objects and then
lexicographically minimize
\[
  (\text{radical depth},\text{radical-node count},\text{leaf count}).
\]
Alternatively, a user might prefer minimum leaf count even at the cost of a slightly
greater nesting depth. Either policy is defensible; silently substituting minimal-
polynomial degree for both is not.

\subsection{Forced multipliers are not necessarily the only multipliers}
A list supplied as the second argument to \code{denest11} replaces the initial
multiplier list and changes its batch flag. It does not disable every path that
calls \code{newmultipliers}.

The worked example supplies a concrete control-flow illustration. With
\code{{1, Sqrt[5]}}, the first trial has degree four and no smaller GCD; the second
has degree eight. While no worse degree is yet known, that worse trial counts as
progress. The subsequent history logic can schedule new discriminant-based
multipliers. Thus an option called ``forced multipliers'' is better described as
``initial multipliers'' unless an additional option explicitly forbids generation.

\subsection{Numerical tests should not be blamed for the wrong reason}
The real-part sign test uses 100-digit numerical evaluation; the complex grouping
rule uses 10-digit numerical evaluation. A fixed precision is not an exact sign or
zero certificate for an arbitrarily ill-conditioned algebraic expression.
Nevertheless, the sign-normalization loop negates both a factor and an accumulated
sign, preserving their product regardless of which branch of that test executes.
Similarly, grouping existing factors into a marked product does not itself change
their mathematical product.

The immediate risks here are representation sensitivity, inconsistent search paths,
and poor behavior near cancellation. The decisive demonstrated wrong-value bug is
the later principal-root reconstruction, not an unsupported claim that every
numerical sign misclassification changes the product's value.

\subsection{F12 and other maintenance issues}
\code{extrapower} is computed from inner nesting levels and printed, but it is never
used to construct a trial, alter $q$, or improve a result. The diagnostic ``extra
root'' should not be taken as evidence that an extra-root algorithm is implemented.
\code{maxprimes = 14} is likewise unused. Other local names are unused or used only
for diagnostics, and obsolete alternatives remain commented out.

\code{fullfactor} stops its local refinement loop after five passes whether or not
\code{done} is true. This is an arbitrary normalization limit, not a documented
mathematical bound. A limit can be a legitimate safeguard, but reaching it should
be explicit and should preserve a well-defined representation rather than silently
suggest completion.

The helper \code{combine} removes unit-base entries only inside a loop ending at
\code{Length[base]-1}; a final or lone unit entry can survive that cleanup. This is
primarily an inconsistency in canonicalization, not on its own a proved wrong-value
case. The subsequent recursive \code{Position} and \code{Delete} operations also
rely on the private list encoding not containing look-alike data.

Finally, undocumented positional records, a shared string tag for
\code{Catch}/\code{Throw}, diagnostic printing embedded in predicates, and global
helper definitions make individual components difficult to test independently.
These are engineering shortcomings distinct from the mathematical counterexamples.

\section{The denominator rationalizer is mathematically sound}
\label{sec:rationalizer}
This helper deserves a positive assessment, subject to its domain assumptions.
For a nonzero algebraic denominator $d$, let
\[
 P(X)=a_nX^n+\cdots+a_1X+a_0
\]
be its minimal polynomial over $\Q$. Because $d\ne0$, one has $a_0\ne0$.
Polynomial division gives
\[
 P(X)=(X-d)Q(X),
\]
with coefficients of $Q$ in $\Q(d)$. Evaluating at $X=0$ gives
\[
 a_0=-dQ(0),
 \qquad \frac1d=-\frac{Q(0)}{a_0}.
\]
That is exactly the formula implemented by
\begin{lstlisting}
PolynomialQuotient[minpoly[x], x - denom, x] /. x -> 0
\end{lstlisting}
followed by division by \code{-minpoly[0]} and multiplication by the numerator.

Equivalently, the annihilating relation $P(d)=0$ yields
\begin{equation}
 \frac1d=-\frac{a_1+a_2d+\cdots+a_nd^{n-1}}{a_0}.
 \label{eq:inverse}
\end{equation}
The stored variable \code{conjugate} is a polynomial, not the ordinary complex
conjugate of $d$; the name is misleading, but the formula is valid.

For $d=\sqrt2+\sqrt3$, the polynomial is $X^4-10X^2+1$, so
\[
 \frac1d=10d-d^3=\sqrt3-\sqrt2.
\]
For $d=1+i\sqrt2$, the polynomial is $X^2-2X+3$, giving
\[
 \frac1d=\frac{2-d}{3}=\frac{1-i\sqrt2}{3}.
\]
Both identities, and a rational-denominator case, are checked independently in the
companion script.

The missing engineering pieces are explicit rejection of zero or unsupported
denominators, checked polynomial operations, and a policy about expression growth.
Also, a sum of rational expressions need not already have the desired common
syntactic denominator. Calling \code{Together} before denominator extraction is a
normalization choice; without it the helper can preserve the value while failing
to rationalize every denominator hidden in separate summands.

\section{A conservative repair and redesign}
\label{sec:repair}

\subsection{First establish a correctness boundary}
The first milestone should not be a more aggressive multiplier generator. It should
be a kernel with the following invariant:
\begin{quote}
Every returned candidate has been proved equal to the original exact algebraic
input; search heuristics can change which candidates are found, but cannot change
what counts as an acceptable answer.
\end{quote}
Concretely, retain the original input, remove branch-unsafe preprocessing, reject
invalid multipliers, enumerate the existing root-of-unity orbit, and compare each
candidate with that original value using exact algebraic arithmetic. If no candidate
is certified, keep the original. At the top-level boundary, certify the reconstructed
whole expression again whenever its domain permits exact comparison.

For symbolic host expressions, exact algebraic comparison of the whole expression
is not available. There the rewrite engine must preserve the host structure and
replace only validated numeric subexpressions, or establish identities under
explicit assumptions. It must not silently treat a symbolic identity as a numeric
zero-testing problem.

\subsection{Give the result an explicit type}
A trial should return a record such as
\begin{lstlisting}
<|"Status" -> "Candidate",
  "Original" -> original,
  "Multiplier" -> multiplier,
  "MinimalPolynomial" -> polynomial,
  "RationalDegree" -> degree,
  "RelativeGCDDegree" -> gcdDegree,
  "Candidates" -> verifiedCandidates|>
\end{lstlisting}
Other statuses should distinguish absence of a candidate, no GCD-degree reduction,
invalid input, failed computation, and an exceeded budget. A search can then record
a worsening degree without
pretending that this is a successful denesting, and it can return its best certified
candidate when a resource budget is reached.

\subsection{A replacement candidate kernel}
The companion \code{safe_candidate_kernel.wl} supplies a small, conservative
implementation of the mathematically justified trial. It does not reproduce the
attachment's scheduler or claim completeness. Its important differences are that
$q$ is explicit and validated, the original value is never discarded, a nonzero
exact multiplier is required, and an empty candidate set is an ordinary result.

The central extraction and verification are:
\begin{lstlisting}
rho = original^q;
beta = (multiplier rho)^(1/q);
polynomial = MinimalPolynomial[beta, x];
g = PolynomialGCD[polynomial,
                  x^q - multiplier rho,
                  Extension -> Automatic];
roots = x /. Solve[g == 0, x];
candidates = Flatten[Outer[Times,
   roots/multiplier^(1/q),
   Exp[2 Pi I Range[0, q - 1]/q]]];
verified = Select[candidates,
   ExactAlgebraicEqualQ[#, original] &];
\end{lstlisting}
The complete file adds guards, return-type checks, exact duplicate elimination, and
a time-budget wrapper. It is a proposed implementation, not a target-kernel-tested
patch. Its algebraic construction is the one proved in Section~\ref{sec:algebra}.

\subsection{Separate equivalence testing from representation scoring}
Do not apply \code{RootReduce} to the retained representation merely because it is
used to prove equivalence. The verification representation and the preferred
output representation have different purposes. An exact \code{Root} object may
be excellent for equality but undesirable as the visible result of a radical
denester.

For root-free numeric inputs, a defensible cost tuple is
\[
 \bigl(\#\text{implicit algebraic objects},\ \depthrad,
        \#\text{rational-power nodes},\ \text{leaf count}\bigr).
\]
Count a rational-power node only when its exponent has denominator greater than
one. Treat \code{Root} and \code{AlgebraicNumber} explicitly; do not accidentally
reward them with a radical depth of zero. Use a strict improvement rule and retain
the best certified answer found so far.

A single traversal can compute depth and other structural statistics. The current
combination of annotated nested lists, repeated rewrites, pruning, and positional
extraction is more complicated than necessary and exposes more representation-
dependent corner cases.

\subsection{Rebuild the scheduler around budgets and a visited set}
The search layer should own a priority queue of untried multipliers, a set of
canonical exact values already attempted, and a cache of trial results. Generated
batches should be lazy streams. Enforce a total trial budget, a maximum candidate
representation size, and time and memory budgets around expensive algebraic
operations. A hard budget must be checked before materializing a cartesian product,
not after it.

Use distinct options for an initial multiplier list and an exclusive multiplier
list. Make heuristic prime pruning optional, record discarded composite residuals,
and expose the reason each candidate was proposed. These controls make a missed
denesting diagnosable rather than turning a heuristic failure into a mysterious
``no progress'' printout.

\subsection{A cheap exact fast path for quadratic surds}
Before invoking the general search, recognize a common special case:
\[
 \sqrt{a+b\sqrt d},\qquad a,b,d\in\Q,\quad d>0.
\]
Suppose $a>0$ and
\[
 \Delta=a^2-b^2d=\delta^2\ge0,\qquad \delta\in\Q_{\ge0}.
\]
Set
\[
 u=\frac{a+\delta}{2},\qquad v=\frac{a-\delta}{2}.
\]
Then $u,v\ge0$, $u+v=a$, and $4uv=b^2d$. Consequently, for $b\ne0$,
\begin{equation}
 \sqrt{a+b\sqrt d}=\sqrt u+\operatorname{sgn}(b)\sqrt v.
 \label{eq:quadfast}
\end{equation}
Indeed the right side squares to the radicand and is nonnegative because $u\ge v$.
The case $b=0$ is simply $\sqrt a$. This derivation also identifies the positivity
and branch hypotheses that a general symbolic rewrite must not omit.

For $a=5$, $b=2$, and $d=6$, one has $\delta=1$, $u=3$, $v=2$ and immediately obtains
\eqref{eq:classic}. Such a fast path avoids minimal polynomials and discriminant
factorization for an important family while retaining exact verification.

\section{Testing strategy and reproducibility}
\label{sec:tests}

\subsection{What was actually executed}
The supplied \code{independent_checks.py} ran successfully and recorded 26 checks.
These cover the classical denesting identity; its minimal polynomial and
discriminant; degrees and GCDs for five representative multipliers; the negative,
three-halves, and merged-product branch constructions; the denominator inverse
formula; and small Python models of the comparator, stale-length, and multiplier-
count defects.

The exact Python and SymPy versions, source hash, symbolic results, and per-check
statuses are in \code{verification_results.json}. The human-readable output is in
\code{verification_results.txt}. These files do not assert that \code{Strad},
\code{Factorc}, or the complete Wolfram search ran successfully or failed in a
particular observed way.

\subsection{Wolfram regressions supplied for execution}
The companion test file contains assertions for correctness, not assertions that
encode the buggy behavior as desirable. Some are therefore expected to fail against
the original code. Important cases are:
\small
\begin{longtable}{@{}p{0.31\linewidth}p{0.59\linewidth}@{}}
\toprule Test family & Required property\\\midrule
\endfirsthead
\toprule Test family & Required property\\\midrule
\endhead
Positive classical radical & The answer equals $\sqrt2+\sqrt3$.\\
Negative direct core input & The sign of $-\sqrt{5+2\sqrt6}$ is preserved.\\
Complex exponent $3/2$ & The answer equals $-5+i\sqrt2$, not its negative.\\
Bundled complex product & Both the direct core and wrapper preserve the product in \eqref{eq:product}.\\
Symbolic factorization & Substitution after normalization does not change the value of $\sqrt{-x-y}$.\\
Zero multiplier & The trial rejects invalid input or conservatively preserves the original value; it does not return a malformed ``answer.''\\
Prime ordering and pruning & The comparator uses both records; deleting primes cannot invalidate a subsequent slice.\\
Global state & Loading the source does not permanently unprotect \code{Print}.\\
All-level operation & The result remains equal to the original nested input under both modes.\\
Replacement candidate kernel & Certified candidates preserve all tested branches and invalid multipliers return failures.\\
\bottomrule
\end{longtable}
\normalsize
The runner records the actual kernel version and uses finite evaluation timeouts.
It restores the incoming attributes of \code{Print} after loading and testing the
source. The package remains an audit bundle, not a claim that the replacement
kernel has completed platform-specific validation.

\subsection{Properties beyond individual examples}
A mature test suite should enforce semantic preservation on generated exact
algebraic inputs, stability under multiplication by roots of unity, and consistency
between equivalent input representations. Tests should compare exact values rather
than one preferred textual form.

For accepted rewrites, verify both equivalence and strict cost improvement. For
unchanged returns, verify the status says only what is known: no improvement found
within the selected strategy and budget. Exercise cancellation near the real and
imaginary axes, rational exponents with numerator greater than one, negative
exponents, mixed root orders, rational denominators, explicit \code{Root} inputs,
invalid forced multipliers, duplicate multipliers, and large prime gaps.

The raw original source, its hash, the independent checker, and the proposed
regressions are included so that a future kernel run can distinguish a source
change, a language-version change, and a mistaken analysis.

\section{Final assessment}
The attachment contains a useful experimental idea: a multiplier can lower the
rational degree of an auxiliary radical, allowing a polynomial GCD over an algebraic
coefficient field to expose a simpler representation. The classical example shows
this mechanism exactly, and the denominator rationalizer is supported by a clean
algebraic proof.

The principal weakness is not merely that the search is heuristic. It is that
normalization, branch selection, proposal generation, and correctness verification
are entangled. A heuristic can be incomplete while remaining perfectly reliable;
this implementation does not maintain that separation. Its wrong-target filter
can approve a sign-reversed number, and its factorization rules can change branches
before the search begins.

The practical repair order is therefore: preserve and certify the original value;
remove or guard unsafe power transformations; fix the comparator, stale slice, and
empty-result paths; define explicit input and output types; then bound and improve
the multiplier search. After that foundation is in place, stronger heuristics and
specialized denesting algorithms can be added without jeopardizing correctness.

\appendix
\section{Companion files and build instructions}
The archive is self-contained. It includes the article source and PDF, the unmodified
Wolfram source, the independent checker and its executed results, a conservative
replacement candidate kernel, Wolfram regression tests and a runner, and a README.

Compile this article from the archive directory with:
\begin{lstlisting}[language={}]
pdflatex -interaction=nonstopmode -halt-on-error review.tex
pdflatex -interaction=nonstopmode -halt-on-error review.tex
\end{lstlisting}
Run the independent checks with a Python installation containing SymPy:
\begin{lstlisting}[language={}]
python independent_checks.py
\end{lstlisting}
Run the Wolfram regressions in a kernel environment with:
\begin{lstlisting}[language={}]
wolframscript -file run_tests.wl
\end{lstlisting}
The tests intentionally express the desired correctness properties, so the original
source is expected to fail several of them. No files containing font binaries are
included or required beyond the fonts supplied by a standard suitable TeX
installation.

\section{Reference implementation of one certified candidate trial}
\label{app:kernel}
This listing is a proposed building block, not a complete replacement for
\code{DenestRadicals3}. It is reproduced from \code{safe_candidate_kernel.wl}.
Its exact-algebraic guards are intentionally conservative: an input whose domain
membership cannot be established is rejected rather than guessed.
\lstinputlisting{safe_candidate_kernel.wl}

\clearpage
\section{The complete original source, with line numbers}
\label{app:source}
The following listing is read directly from the byte-preserved \code{source.wl}.
Long physical lines may wrap visually, but the printed line numbers refer to the
original physical lines. No repairs have been applied to this listing.
\lstinputlisting[style=wl,basicstyle=\ttfamily\scriptsize,numbers=left,
 numberstyle=\tiny\color{muted},stepnumber=1,numbersep=7pt,
 xleftmargin=2.1em,backgroundcolor=\color{white},frame=none]{source.wl}

\clearpage
\begin{thebibliography}{99}
\bibitem{source} Supplied source, \emph{Pasted text(20260905-171712).txt}.
Unmodified copy distributed as \code{source.wl}; 585 physical lines. Source-line
references throughout this article identify the analyzed definitions directly.

\bibitem{doc-power} Wolfram Research, \emph{Power}, Wolfram Language documentation.
\url{https://reference.wolfram.com/language/ref/Power.html}.
Accessed September 5, 2026.
\bibitem{doc-minpoly} Wolfram Research, \emph{MinimalPolynomial}.
\url{https://reference.wolfram.com/language/ref/MinimalPolynomial.html}.
Accessed September 5, 2026.
\bibitem{doc-gcd} Wolfram Research, \emph{PolynomialGCD}.
\url{https://reference.wolfram.com/language/ref/PolynomialGCD.html}.
Accessed September 5, 2026.
\bibitem{doc-possiblezero} Wolfram Research, \emph{PossibleZeroQ}.
\url{https://reference.wolfram.com/language/ref/PossibleZeroQ.html}.
Accessed September 5, 2026.
\bibitem{doc-rootreduce} Wolfram Research, \emph{RootReduce}.
\url{https://reference.wolfram.com/language/ref/RootReduce.html}.
Accessed September 5, 2026.
\bibitem{doc-slot} Wolfram Research, \emph{Slot}.
\url{https://reference.wolfram.com/language/ref/Slot.html}.
Accessed September 5, 2026.
\bibitem{doc-sort} Wolfram Research, \emph{Sort}.
\url{https://reference.wolfram.com/language/ref/Sort.html}.
Accessed September 5, 2026.
\bibitem{doc-factorinteger} Wolfram Research, \emph{FactorInteger}.
\url{https://reference.wolfram.com/language/ref/FactorInteger.html}.
Accessed September 5, 2026.
\bibitem{doc-block} Wolfram Research, \emph{Block}.
\url{https://reference.wolfram.com/language/ref/Block.html}.
Accessed September 5, 2026.
\bibitem{doc-apply} Wolfram Research, \emph{Apply}.
\url{https://reference.wolfram.com/language/ref/Apply.html}.
Accessed September 5, 2026.
\bibitem{doc-memberq} Wolfram Research, \emph{MemberQ}.
\url{https://reference.wolfram.com/language/ref/MemberQ.html}.
Accessed September 5, 2026.
\bibitem{doc-complexexpand} Wolfram Research, \emph{ComplexExpand}.
\url{https://reference.wolfram.com/language/ref/ComplexExpand.html}.
Accessed September 5, 2026.
\bibitem{doc-powerexpand} Wolfram Research, \emph{PowerExpand}.
\url{https://reference.wolfram.com/language/ref/PowerExpand.html}.
Accessed September 5, 2026.
\bibitem{doc-algebraics} Wolfram Research, \emph{Algebraics}.
\url{https://reference.wolfram.com/language/ref/Algebraics.html}.
Accessed September 5, 2026.
\end{thebibliography}
\end{document}
